\documentclass[10pt,a4paper]{amsart}
\usepackage[margin=19mm]{geometry}
\usepackage{amsmath,amssymb,mathtools}
\usepackage[scr=rsfs,cal=euler]{mathalfa}
\usepackage{xfrac}
\usepackage[unicode,hidelinks,bookmarks=false]{hyperref}
\newcommand{\ie}{i.e.\ }
\newcommand{\cf}{cf.\ }
\newcommand{\resp}{respectively\ }
\newcommand{\Subsection}{\subsection}
\providecommand{\nobreakdash}{\nobreak\hbox{-}}
\setlength{\emergencystretch}{2em}
\multlinegap0pt
\usepackage{bm}
\usepackage{bbm}
\usepackage{dsfont}
\usepackage{xspace}
\usepackage{cancel}
\usepackage{mathtools}
\usepackage{amsthm}
\makeatletter
\@ifundefined{theorem}{\newtheorem{theorem}{Theorem}}{}
\@ifundefined{lemma}{\newtheorem{lemma}[theorem]{Lemma}}{}
\makeatother
\title[Coloring t-perfect graphs with fewer colors]{Coloring t-perfect graphs with fewer colors}
\author{Matija Novakovi\'c, Stefan Weltge}
\date{Source: arXiv:2607.05961v1 (CC BY 4.0)}
\begin{document}
\maketitle

\noindent\emph{Source and licence.} Coloring t-perfect graphs with fewer colors. Version arXiv:2607.05961v1 (CC BY 4.0), distributed under the Creative Commons Attribution 4.0 International licence, as recorded in the arXiv licence metadata for this exact version and shown on the abstract page https://arxiv.org/abs/2607.05961v1. Full rights notice: licenses/arxiv-2607.05961-CC-BY-4.0.txt.\par\smallskip

\noindent\textit{Editorial scope.} Counted unit: the Lemma whose source label is \texttt{lemma1:beyond\_leveling} in the article (source0.tex lines 424--435, source label \texttt{lemma1:beyond\_leveling}), delivered together with the article's own complete original proof of it (source0.tex lines 438--459, one \texttt{proof} environment of 22 lines). No mathematics is written, repaired or re-derived: statement and proof are the article's own text line for line. Required context retained uncounted: none. Every definition, display and label of those lines is retained unchanged. Omitted from this package and not redistributed: 1--423, 436--437, 460--736. Nothing inside a retained excerpt is omitted or altered: every retained body is delivered exactly as the article's lines are written, so no omission receipt of any kind accompanies this package. Citations: the retained excerpts carry no citation command at all, so no bibliography entry of the article is transcribed and no reference is redistributed. Reference adaptation: none was needed and none was made; every reference inside the delivered unit resolves inside it, so no unresolved reference or citation is produced. Printed slips kept literal: no printed word, symbol, hypothesis or number of the counted statement or of its proof was altered. Layout and class: the article's own class is \texttt{article} with packages including inputenc, fullpage, amsmath, amsthm, amssymb, mathtools, tikz, hyperref, cleveref, todonotes, xspace; the wrapper is the lane's pinned \texttt{amsart} editorial wrapper at 10pt on A4, which restates the theorem-like environments the retained text needs and adds the article's own macro definitions that the retained text needs (none), with extra wrapper packages (bm, bbm, dsfont, xspace, cancel, mathtools). The delivered numbering is the unit's own. Assets: no retained span contains a figure environment, a graphics-inclusion command, a PDF-page inclusion, a table or a TikZ picture; no image file is copied into this package, the package declares no asset dependency and no image file is redistributed with it. Licence: version arXiv:2607.05961v1 of this article is distributed under the Creative Commons Attribution 4.0 International licence, as recorded in the arXiv licence metadata for this exact version and shown on the abstract page https://arxiv.org/abs/2607.05961v1. A full rights notice is delivered at licenses/arxiv-2607.05961-CC-BY-4.0.txt. Characters: every retained excerpt is pure ASCII, so the delivered engine, XeLaTeX with its default Latin Modern text font, reports no missing character.
\begin{lemma}\label{lemma1:beyond_leveling}
    Let $c \geq 1$ and let $G$ be a graph with $\chi(G) \geq c + 8$ and $V(G) = C_0 \cup C_1 \cup C_2$.
    Furthermore, let \((C_0^t)_{t\geq 1}\) be a partition of \(C_0\) such that the ends of every edge of \(G[C_0]\) lie in parts whose indices differ by at most one.

    If $(W_1, \dots, W_n)$ is a stable grading of $V(G)$, there exists a subset $X$ of $V(G)$ and an edge $uv$ belonging to one of $G[C_1], G[C_2]$, or $G[C_0^t]$, for some $t \geq 1$, such that
    \begin{itemize}
        \item $G[X]$ is connected,
        \item $\chi(X) \geq c$,
        \item $u$ and $v$ are earlier than every vertex in $X$, and
        \item at least one of $u$ and $v$ has a neighbor in $X$.
    \end{itemize}
        \end{lemma}

\begin{proof}
    Let $A$ be the set of vertices in $G$ that are left-active with respect to $C_1, C_2,$ or $C_0^t$ for some $t \geq 1$. 

    \textbf{Case 1:} $\chi(A) \geq c$.
    Let $D$ be a connected component of $G[A]$ with maximum chromatic number. Note that we have $\chi(D) = \chi(A) \geq c$. Then there is a minimal integer $i$ such that $1 \leq i \leq n$ and $W_i \cap V(D)$ is non-empty. Let $w$ be a vertex in $W_i \cap V(D)$. Since $w$ is left-active, there is an edge $uv$ of $G[C_1], G[C_2]$ or $G[C_0^t]$, for some $t \geq 1$, such that $u$ and $v$ are earlier than $w$ and $w$ is adjacent to one of $u$ and $v$. Since $i$ was chosen to be minimal, $u$ and $v$ are earlier than every vertex in $D$. Thus $X = V(D)$ satisfies the required conditions.

    \textbf{Case 2:} $\chi(A) \leq c - 1$.
    Let $B_h = C_h \setminus A$ for $h \in \{0, 1, 2\}$. We perform a case distinction depending on which $B_h$ has a sufficiently large chromatic number.
    Since 
    \[\chi(B_0 \cup B_1 \cup B_2) \geq \chi(G) - \chi(A) \geq c + 8 - (c -1) = 9,\]
    we either have $\chi(B_0) \geq 5$ or $\chi(B_1 \cup B_2) \geq 5$. In the second case, one of $B_1$ or $B_2$ must have chromatic number at least $3$.

    \textbf{Case 2.1:} $\chi(B_0) \geq 5$.
    Since distinct same-parity parts among the \(C_0^t\)'s are anticomplete in \(G[C_0]\), there is an integer $r \geq 1$ with $\chi(B_0 \cap C_0^{r}) \geq \left\lceil \chi(B_0) / 2 \right\rceil \geq 3$.
    Hence \(G[B_0 \cap C_0^r]\) contains an odd cycle \(D\) as a subgraph. Since each \(W_i\) is stable, the endpoints of every edge of \(D\) lie in different parts of the stable grading. Orient each edge from the earlier endpoint to the later endpoint, and denote the resulting orientation by \(\vec D\).

    Since no vertex in \(D\) is left-active with respect to \(C_0^r\), there is no directed path of length \(2\) in \(\vec D\). Indeed, if \(u \to v \to w\) is such a path, then \(uv\) is an edge of \(G[C_0^r]\), both \(u\) and \(v\) are earlier than \(w\), and \(w\) is adjacent to \(v\), contradicting \(w\in B_0\).
    Hence every vertex of \(\vec D\) is either a sink or a source. Thus the sources and sinks form a bipartition of \(D\), contradicting that \(D\) is an odd cycle. 

    \textbf{Case 2.2:} $\chi(B_h) \geq 3$ for some $h \in \{1, 2\}$.
    The same orientation argument applied to an odd cycle in \(G[B_h]\) gives a contradiction, since no vertex of \(B_h\) is left-active with respect to \(C_h\).
\end{proof}

\end{document}
